% introduccion
\vspace{1em}
Al momento de establecer una comunicación entre dos o más computadoras, se utilizan los llamados Medios de Acceso Compartido. Estos son canales de transmisión de información donde varios dispositivos se ponen de acuerdo sobre un conjunto de reglas y protocolos para poder intercambiar datos.

\vspace{1em}
A su vez, los datos transmitidos son divididos en paquetes más pequeños, también llamados tramas, para su transportación. Cada uno de ellos contiene información sobre los protocolos que encapsula, su origen, y su destino particular, en caso de un paquete de tipo \textit{Unicast}, o pueden indicar que fueron enviados a todos los nodos de una red, para los paquetes de tipo \textit{Broadcast}. Estas tramas pueden ser capturadas y estudiadas por cualquier dispositivo conectado al medio. 

\vspace{1em}
A lo largo de este trabajo analizaremos las características de tres redes domésticas particulares. En especial, capturaremos los diversos paquetes de datos que las circulan, analizaremos los protocolos que utilizan las distintas capas de red, y buscaremos tendencias y patrones entre las tramas.
